Much of what a language model does inside a production system is not writing but deciding: which tool to call next, whether a cached answer can be reused, which model should answer.
A new class of non-generative ``decision models'' returns a typed choice with a probability per option instead of text and is marketed as a far cheaper replacement for an LLM at these points.
We test that claim on one harness with interchangeable deciders (a frontier LLM, small hosted and self-hosted LLMs, a proprietary decision model, an open-weights reproduction of it, and trained classifiers) at three decision points, measuring quality, cost, latency and GPU energy.
The decision model does not pay off for agent tool choice, where letting the generator decide inside its own call is cheapest at equal task reward.
It pays off as a semantic-cache guard, which decides whether a stored answer can be served for a new query: on math questions whose near-identical variants may change the answer, it decides correctly 93.3\% of the time, against 81.7\% for a classifier trained on that dataset and 66.1\% for a similarity threshold, at about 2\% of the frontier model's cost per decision.
For model routing, a small classifier trained on the cheap model's record beats every general decider.
Across the three decision points, a separate decider was worth its call only when the decision avoids a generation and the evidence it needs is in the input.
